\subsection{Introducción del trabajo}

El presente trabajo de investigación se desarrolla en el marco de la asignatura
Sistemas de Comunicaciones de la Universidad Tecnológica Nacional. Su propósito
es analizar ROS (\textit{Robot Operating System}) desde la perspectiva de los
sistemas de comunicaciones, considerando los mecanismos de intercambio de
información que emplea y los estándares sobre los que se sustenta.

Los sistemas robóticos actuales constituyen, en esencia, sistemas distribuidos.
Sensores, actuadores, unidades de procesamiento y estaciones de supervisión
operan como procesos independientes, frecuentemente alojados en distintos
equipos, que deben intercambiar datos de manera continua a través de una red.
En este contexto, la arquitectura de comunicaciones deja de ser un aspecto
secundario y pasa a condicionar directamente la latencia, la confiabilidad y la
escalabilidad del sistema en su conjunto.

ROS se ha consolidado como el \textit{framework} de referencia para el
desarrollo de software robótico, tanto en el ámbito académico como en el
industrial. Su evolución resulta particularmente relevante para el estudio de
las comunicaciones: la primera versión (ROS~1) se basa en una arquitectura
híbrida que depende de un nodo maestro central, mientras que ROS~2 adopta un
esquema descentralizado sustentado en el estándar DDS (\textit{Data
Distribution Service}), el protocolo RTPS (\textit{Real-Time Publish-Subscribe})
y el transporte sobre UDP/IP. Este trabajo aborda dicha transición, analizando
sus fundamentos técnicos, la configuración de la calidad de servicio (QoS) y
las implementaciones de \textit{middleware} disponibles.

\subsection{Objetivo y alcance}

El objetivo general del trabajo es analizar el estado del arte de ROS, con foco
en la evolución de su arquitectura de comunicaciones desde ROS~1 hasta ROS~2 y
en los estándares y protocolos que la sustentan.

Para ello se plantean los siguientes objetivos específicos:
\begin{itemize}
    \item Describir los conceptos arquitectónicos de ROS (nodos y \textit{topics})
    y sus patrones de comunicación (\textit{services} y \textit{actions}).
    \item Analizar la arquitectura híbrida de ROS~1 y las limitaciones derivadas
    de su dependencia del nodo maestro (\textit{roscore}) como punto único de
    falla.
    \item Describir el modelo descentralizado de ROS~2 y la pila de
    comunicaciones DDS, RTPS y UDP/IP.
    \item Analizar los mecanismos de calidad de servicio (QoS) y su
    configuración.
    \item Relevar el soporte de ROS~2 para sistemas operativos de tiempo real
    (RTOS) y microcontroladores mediante micro-ROS.
    \item Describir la capa de abstracción de \textit{middleware} (RMW) y
    comparar sus principales implementaciones: Eclipse Cyclone DDS, eProsima
    Fast DDS y RTI Connext DDS.
    \item Presentar el esquema de distribuciones de ROS y su vinculación con
    las versiones de Ubuntu.
    \item Realizar una demostración práctica sobre hardware dedicado (Raspberry
    Pi) que ilustre el intercambio de mensajes entre nodos de ROS~2 a través de
    la red.
\end{itemize}

En cuanto al alcance, el estudio se centra en la capa de \textit{middleware} y
transporte de datos, complementado con una prueba de concepto de carácter
demostrativo. No se abordan en profundidad los algoritmos de
percepción, planificación o control, ni el desarrollo de aplicaciones
robóticas específicas.

\subsection{Metodología y Fuentes de investigación}
La investigación si el trabajo de investigación es del tipo descriptivo y posee
atributos cualitativos o tiene otro foco.

\subsection{Delimitación del estudio}
La delimitación del análisis se da en las áreas por ejemplo de:

- Sector Telecomunicaciones
- Producto: servicio de telefonía móvil

\subsection{Universo Espacio Temporal}
Los años en los cuales se tuvo presente la información para hacer el trabajo.